\documentclass[12pt,a4paper]{article}

% ----------------------------------------------------------------------
% Codificação, fonte e idioma
% ----------------------------------------------------------------------
\usepackage[utf8]{inputenc}
\usepackage[T1]{fontenc}
\usepackage[brazilian]{babel}
\usepackage{lmodern}
\usepackage{microtype}

% ----------------------------------------------------------------------
% Layout
% ----------------------------------------------------------------------
\usepackage[a4paper,top=3cm,bottom=2cm,left=3cm,right=2cm]{geometry}
\setlength{\headheight}{14pt}
\usepackage{setspace}
\onehalfspacing
\usepackage{indentfirst}
\usepackage{graphicx}
\usepackage{booktabs}
\usepackage{longtable}
\usepackage{array}
\usepackage{tabularx}
\usepackage{enumitem}
\usepackage{xcolor}
\usepackage{listings}
\usepackage{caption}
\usepackage{float}
\usepackage{tikz}
\usetikzlibrary{arrows.meta,positioning,shapes.geometric}
\usepackage[hidelinks]{hyperref}
\hypersetup{
  pdftitle={Chatbot ITT 2026.2: propósito, evolução e contribuições},
  pdfauthor={Guilherme de Godoi Monteiro; Henrique Ribeiro Medeiros da Silva; Renilson Conceição de Luna Júnior; Vitor Ferreira dos Santos de Lacerda},
  pdfsubject={Relatório da disciplina MX853},
  pdfkeywords={RAG, chatbot, FAISS, Google Drive, FastAPI, LangGraph}
}
\usepackage{titlesec}
\usepackage{fancyhdr}

\definecolor{ittred}{HTML}{B91C1C}
\definecolor{codebg}{HTML}{F5F5F5}
\definecolor{codecomment}{HTML}{6B7280}

\lstset{
  basicstyle=\ttfamily\small,
  backgroundcolor=\color{codebg},
  commentstyle=\color{codecomment},
  frame=single,
  rulecolor=\color{gray!40},
  breaklines=true,
  columns=fullflexible,
  keepspaces=true,
  showstringspaces=false,
  numbers=left,
  numberstyle=\tiny\color{gray},
  numbersep=8pt,
  xleftmargin=1.5em,
  framexleftmargin=1.2em
}

\pagestyle{fancy}
\fancyhf{}
\fancyhead[L]{\small MX853 -- Projeto Integrador Interdisciplinar de Extensão}
\fancyhead[R]{\small Chatbot ITT 2026.2}
\fancyfoot[C]{\thepage}
\renewcommand{\headrulewidth}{0.4pt}

\setlist[itemize]{topsep=4pt,itemsep=2pt,parsep=0pt}
\setlist[enumerate]{topsep=4pt,itemsep=2pt,parsep=0pt}

\newcommand{\commit}[1]{\texttt{#1}}
\newcommand{\arquivo}[1]{\path{#1}}

% ----------------------------------------------------------------------
% Capa
% ----------------------------------------------------------------------
\begin{document}

\begin{titlepage}
  \thispagestyle{empty}
  \begin{center}
    {\Large Universidade Estadual de Campinas}\\[0.2cm]
    {\large Instituto de Computação}\\[2.0cm]

    {\large MX853 -- Projeto Integrador Interdisciplinar de Extensão}\\[2.5cm]

    {\Huge\bfseries Chatbot ITT 2026.2}\\[0.5cm]
    {\Large Propósito, evolução e contribuições da equipe}\\[2.5cm]

    \begin{tabular}{ll}
      Guilherme de Godoi Monteiro & 237031\\
      Henrique Ribeiro Medeiros da Silva & 243172\\
      Renilson Conceição de Luna Júnior & 247358\\
      Vitor Ferreira dos Santos de Lacerda & 231524\\
    \end{tabular}

    \vfill
    {\large Orientação: Profª. Drª. Juliana Freitag Borin}\\[0.5cm]
    {\large Campinas -- SP}\\[0.2cm]
    {\large Novembro de 2026}
  \end{center}
\end{titlepage}

% ----------------------------------------------------------------------
% Resumo
% ----------------------------------------------------------------------
\section*{Resumo}
\addcontentsline{toc}{section}{Resumo}

Este documento descreve o propósito, a evolução e as contribuições da equipe no
desenvolvimento do chatbot do Instituto Tadao Takahashi (ITT) durante o segundo
semestre de 2026. O projeto parte de um sistema antecessor de perguntas e respostas
sobre documentos institucionais, baseado em \emph{Retrieval-Augmented Generation}
(RAG), e avança para uma aplicação interna com autenticação de contas Google
restrita ao domínio do ITT, ingestão de PDFs, Google Docs e planilhas do Google
Drive, consultas estruturadas a fichas de alunos, transmissão de respostas em
\emph{streaming}, exibição de fontes, testes automatizados e execução com Docker
Compose. O texto está organizado em três eixos: o propósito do projeto, as adições
feitas em relação ao código herdado e a descrição do que foi realizado por cada
participante, com indicação dos commits e arquivos que sustentam cada afirmação.

\noindent\textbf{Palavras-chave:} RAG; chatbot; FAISS; Google Drive; FastAPI;
LangGraph; extensão universitária.

\newpage
\begingroup
\small
\tableofcontents
\endgroup
\newpage
\listoffigures
\listoftables
\newpage

% ======================================================================
\section{Propósito do projeto}
% ======================================================================

\subsection{Contexto e problema}

O Instituto Tadao Takahashi (ITT) é uma organização sem fins lucrativos
sediada em Campinas. Sua atuação envolve projetos educacionais, tecnológicos e
sociais, o que produz uma grande quantidade de documentos: estatuto, contratos,
relatórios, planos de trabalho, materiais de formação e fichas de acompanhamento
de alunos. Esses arquivos ficam distribuídos em pastas compartilhadas do Google
Drive e, em muitos casos, só são localizados por quem conhece previamente a
estrutura das pastas.

A busca manual por informação tem três consequências principais. Primeiro, aumenta
o tempo necessário para responder a perguntas simples da coordenação e dos
mentores. Segundo, dificulta a recuperação de registros históricos, como a última
observação de um mentor sobre um aluno ou a quantidade de cursos concluídos em uma
planilha. Terceiro, expõe dados sensíveis: as fichas de alunos contêm informações
pessoais e de acompanhamento que não devem ser tratadas como texto genérico.

O projeto se insere nesse contexto como uma iniciativa de extensão universitária.
O objetivo não é construir um chatbot genérico, mas uma ferramenta interna que
responda em português com base nos documentos do ITT, indique de onde retirou cada
informação e respeite as restrições de acesso da instituição.

\subsection{Objetivos}

O objetivo geral é disponibilizar um assistente conversacional interno que facilite
a consulta a documentos institucionais e a fichas de alunos, com respostas
fundamentadas e rastreáveis.

Os objetivos específicos são:

\begin{itemize}
  \item ingerir PDFs, Google Docs exportados como PDF e planilhas \texttt{.xlsx}
        armazenados no Google Drive, inclusive em subpastas;
  \item indexar o conteúdo em uma base vetorial FAISS e recuperar os trechos mais
        relevantes para cada pergunta;
  \item gerar respostas em português com um modelo de linguagem, usando somente o
        contexto recuperado e citando as fontes;
  \item autenticar contas Google no backend, restringir o acesso ao domínio do ITT
        e limitar a sincronização de documentos a administradores;
  \item interpretar fichas de alunos de forma estruturada, com filtros por pessoa,
        data, semana, curso, avaliação, frequência e intervalo de células;
  \item transmitir a resposta em \emph{streaming}, permitir interromper a geração e
        exibir as fontes sob demanda;
  \item manter testes automatizados, documentação de execução e um fluxo de
        avaliação de modelos;
  \item empacotar backend e frontend com Docker Compose para reduzir o esforço de
        instalação.
\end{itemize}

\subsection{Público-alvo, escopo e limites}

O público-alvo é a equipe do ITT: coordenação, mentores e administradores que
precisam consultar documentos e acompanhar alunos. A aplicação não é aberta ao
público e depende de login Google.

O escopo atual inclui documentos institucionais em PDF, Google Docs exportados como
PDF e planilhas \texttt{.xlsx}. Também inclui fichas individuais de alunos com abas
de perfil, cursos da Alura, interações com mentores e coordenação, avaliações de
atividades e \emph{scorecard} semanal.

Os principais limites são:

\begin{itemize}
  \item não há permissão por perfil: qualquer pessoa autenticada vê o conteúdo
        indexado;
  \item PDFs digitalizados precisam de OCR (reconhecimento óptico de caracteres) prévio, que não é executado pelo sistema;
  \item arquivos \texttt{.xls} e Google Sheets nativos não são suportados;
  \item fórmulas de planilhas não são recalculadas; o sistema lê os valores salvos;
  \item o cache de embeddings não é podado automaticamente quando um documento sai
        da pasta;
  \item a avaliação de qualidade das respostas ainda depende, em parte, de revisão
        manual.
\end{itemize}

\subsection{Arquitetura e fluxo de uso}

A aplicação é dividida em um frontend React/Vite e um backend FastAPI. O backend
centraliza autenticação, limite de requisições, orquestração do fluxo RAG e acesso
ao índice vetorial. O modelo de linguagem e os embeddings são acessados por meio da
API OpenWebUI da Unicamp, com possibilidade de usar o provedor Gemini como
alternativa legada. O Google Drive é acessado por uma conta de serviço e o login é
validado com o Google Identity.

A Figura~\ref{fig:arquitetura} resume a arquitetura.

\begin{figure}[htbp]
\centering
\resizebox{\textwidth}{!}{%
\begin{tikzpicture}[
  box/.style={draw, rounded corners, align=center, minimum width=3.1cm,
              minimum height=1.05cm, font=\small, fill=white},
  ext/.style={draw, rounded corners, align=center, minimum width=3.1cm,
              minimum height=1.05cm, font=\small, fill=gray!8},
  arr/.style={-{Latex[length=2mm]}, thick}
]
  \node[box] (browser) at (0,0) {Navegador};
  \node[box] (front) at (3.6,0) {Frontend\\React/Vite};
  \node[box] (api) at (7.4,0) {Backend FastAPI\\rotas e validação};

  \node[ext] (identity) at (5.4,-1.7) {Google Identity\\ID token};
  \node[box] (graph) at (9.2,-1.7) {LangGraph\\triagem + RAG};

  \node[box] (faiss) at (7.4,-3.4) {FAISS\\índice + cache};
  \node[box] (llm) at (11.6,-3.4) {OpenWebUI/Unicamp\\Gemma 3 12B + Nomic};

  \node[ext] (drive) at (7.4,-5.1) {Google Drive\\PDF, Docs, XLSX};

  \draw[arr] (browser) -- (front);
  \draw[arr] (front) -- (api);
  \draw[arr] (api) -- (identity);
  \draw[arr] (api) -- (graph);
  \draw[arr] (graph) -- (faiss);
  \draw[arr] (graph) -- (llm);
  \draw[arr] (faiss) -- (drive);
\end{tikzpicture}%
}
\caption{Arquitetura simplificada do chatbot ITT 2026.2.}
\label{fig:arquitetura}
\end{figure}

O fluxo de indexação começa no Google Drive. O backend lista a pasta e suas
subpastas, baixa PDFs e planilhas, exporta Google Docs como PDF e monta uma pasta
temporária. Em seguida, o conteúdo é dividido em trechos, convertido em embeddings
e salvo em um índice FAISS. Se o download ou a leitura falhar, o índice anterior é
preservado.

O fluxo de pergunta começa no navegador. O frontend envia a mensagem e as últimas
mensagens da conversa para o backend. O backend valida o token Google, aplica o
limite de requisições e aciona o grafo LangGraph. O grafo faz uma triagem para
identificar a intenção, busca os trechos no FAISS e gera a resposta com o modelo.
No endpoint de \emph{streaming}, as fontes são enviadas antes dos tokens e o texto
é transmitido em eventos NDJSON.

\subsection{Metodologia e recorte do documento}

O documento foi construído a partir de quatro fontes: o código atual, o histórico
de commits, a documentação existente no repositório e os relatórios do projeto
antecessor. Foram usados comandos como \texttt{git log --all},
\texttt{git diff 087a682..99da614}, \texttt{git show} dos commits relevantes e
leitura dos arquivos em \arquivo{backend/} e \arquivo{frontend/}.

Para evitar ambiguidade, este documento adota como baseline o commit
\commit{087a682}, último commit do projeto antecessor e pai do primeiro commit do
grupo. O primeiro commit do grupo, \commit{804c877}, altera apenas o README; o
primeiro commit com código é \commit{5b3f952}. O estado atual corresponde ao commit
\commit{99da614} da branch \texttt{main}.

A atribuição de contribuições segue a autoria registrada no Git. Commits de merge
são contados separadamente e branches não integradas são identificadas como tal.
Essa escolha não mede esforço individual: um único commit pode concentrar uma
funcionalidade inteira, e o histórico não permite inferir programação em par.
Por segurança, este documento não reproduz credenciais, chaves de API, IDs de
pastas do Drive ou dados de alunos; nomes e RAs aparecem apenas para atribuição
de autoria.

% ======================================================================
\section{Ponto de partida: o projeto antecessor}
% ======================================================================

\subsection{Linhagem e marcos}

O repositório atual é um fork de \texttt{chraujo/chatbot-itt-2026-1}, que por sua
vez deriva de \texttt{programacao-dev/chatbot-itt-}. A Tabela~\ref{tab:marcos}
resume os marcos usados no documento.

\begin{table}[H]
\centering
\small
\begin{tabularx}{\textwidth}{@{}l l l X@{}}
\toprule
Marco & Commit & Data & Significado \\
\midrule
Origem do repositório & \commit{b2db599} & 21/11/2025 &
Commit inicial, ainda do projeto antecessor. \\
Fim do antecessor & \commit{087a682} & 17/07/2026 &
Último commit antes do grupo; código pronto recebido. \\
Primeiro commit do grupo & \commit{804c877} & 23/10/2026 &
Altera apenas o README; primeiro código em \commit{5b3f952}. \\
Estado atual & \commit{99da614} & 27/11/2026 &
Branch \texttt{main} após o merge do PR \#3. \\
\bottomrule
\end{tabularx}
\caption{ Marcos temporais adotados no documento.}
\label{tab:marcos}
\end{table}

O README do projeto antecessor registra Chris Araújo, Richard Matias e Tallyta
Duarte como participantes do primeiro semestre de 2026. O histórico também mostra
contribuições de Antonio, Rodrigo e Marconi em etapas anteriores, incluindo a
primeira integração com o Google Drive e o botão de atualização de documentos.

\subsection{O que veio pronto}

O código disponível no baseline já implementava um fluxo RAG funcional. A
Tabela~\ref{tab:herdado} lista as principais funcionalidades herdadas.

\begin{table}[H]
\centering
\small
\begin{tabularx}{\textwidth}{@{}p{4.0cm} X p{3.3cm}@{}}
\toprule
Funcionalidade & Descrição no baseline & Evidência \\
\midrule
RAG com Gemini & Triagem e geração de resposta com LangGraph, LangChain e
Google Gemini. & \texttt{chains.py}\newline \texttt{graph.py} \\
Índice FAISS & Indexação de PDFs e busca por similaridade. &
\arquivo{vectorDB.py} \\
Google Drive & Download de PDFs de uma pasta e recriação do índice. &
\arquivo{google_drive.py} \\
Auto-sync & Sincronização automática ao iniciar o servidor. &
\arquivo{api.py} \\
Botão ``Atualizar Docs'' & Sincronização manual pela interface. &
\arquivo{Chat.tsx}, \arquivo{admin.py} \\
Rate limit e retry & Limite de 5 requisições por minuto e retry no cliente
Gemini. & \commit{917d497} \\
Rastreamento de fontes & Campos \texttt{source}, \texttt{page} e
\texttt{snippet} na resposta. & \commit{d7a43c0} \\
Login Google & Botão de login no frontend, sem validação no backend. &
\arquivo{Login.tsx} \\
Tema e timestamps & Modo claro/escuro e horário das mensagens. &
\commit{3410322}, \commit{087a682} \\
Docker do backend & Imagem do backend com Python e dependências. &
\arquivo{backend/Dockerfile} \\
\bottomrule
\end{tabularx}
\caption{Funcionalidades já presentes no baseline.}
\label{tab:herdado}
\end{table}

Os relatórios em \arquivo{2026.1/} documentam duas dessas frentes: a otimização do
rate limit e do retry no Gemini e a implementação do rastreamento de fontes. Eles
foram usados como fonte para descrever o ponto de partida, mas não devem ser
confundidos com as contribuições do grupo atual.

\subsection{Lacunas identificadas}

Apesar de funcional, o baseline tinha limitações importantes para o uso interno e
para o tratamento de dados de alunos. A Tabela~\ref{tab:lacunas} resume as
principais.

\begin{table}[H]
\centering
\small
\begin{tabularx}{\textwidth}{@{}p{4.6cm} X@{}}
\toprule
Lacuna & Consequência \\
\midrule
Token Google não validado no backend & Qualquer token poderia ser aceito;
não havia garantia de que o usuário pertencia ao ITT. \\
Sem restrição de domínio ou administradores & Qualquer pessoa logada poderia
acionar a sincronização de documentos. \\
Apenas PDFs & Planilhas e fichas de alunos não eram indexadas nem consultadas
de forma estruturada. \\
Dependência exclusiva do Gemini & Custo, cota e disponibilidade ficavam
condicionados a um único provedor. \\
Sem streaming & O usuário esperava a resposta completa antes de ver qualquer
texto. \\
Sem testes automatizados & Alterações no fluxo RAG, no Drive ou na autenticação
podiam introduzir regressões silenciosas. \\
Sem Docker Compose & A instalação dependia de passos manuais no backend e no
frontend. \\
Sem avaliação de modelos & Não havia forma padronizada de comparar respostas,
fontes e tempos. \\
Histórico global no navegador & As conversas ficavam no \texttt{localStorage} sem
separação por usuário. \\
Sem desambiguação de nomes & Havia risco de misturar fichas de pessoas com nomes
parecidos. \\
\bottomrule
\end{tabularx}
\caption{Lacunas do baseline que motivaram as adições do grupo.}
\label{tab:lacunas}
\end{table}

% ======================================================================
\section{Adições feitas pelo grupo}
% ======================================================================

Esta seção descreve as adições em relação ao baseline. A organização é temática,
e não cronológica, para facilitar a leitura. Cada subseção indica o que existia
antes, o que mudou e quais commits ou arquivos sustentam a descrição. No total,
o diff entre \commit{087a682} e \commit{99da614} altera 49 arquivos, com 4.521
linhas adicionadas e 2.686 removidas.

\subsection{Provedor de LLM e embeddings da Unicamp}

O baseline usava exclusivamente o Google Gemini. O grupo introduziu um provedor
alternativo baseado na API OpenWebUI da Unicamp, com o modelo \texttt{gemma3:12b}
para chat e \texttt{nomic-embed-text:latest} para embeddings.

O arquivo \arquivo{backend/src/services/openwebui.py} implementa um
\texttt{BaseChatModel} e um \texttt{Embeddings} do LangChain sobre os endpoints
\path{/api/chat/completions} e \path{/ollama/api/embed}. A implementação
suporta chamadas síncronas, \emph{streaming} síncrono e \emph{streaming}
assíncrono, necessário para o endpoint de chat em tempo real. Os textos de
indexação recebem o prefixo \path{search_document:} e as consultas recebem
\path{search_query:}, conforme a recomendação do modelo Nomic.

A configuração passou a aceitar \path{LLM_PROVIDER=openwebui|google},
\path{OPENWEBUI_BASE_URL}, \path{OPENWEBUI_API_KEY},
\path{OPENWEBUI_EMBEDDING_MODEL}, \path{LLM_MODEL} e
\path{LLM_TIMEOUT}. O provedor Gemini continua disponível como alternativa
legada.

Para lidar com a cota do servidor, a indexação usa retry exponencial em caso de
HTTP 429, com lotes de 32 trechos e mensagens de progresso. As perguntas do chat
não entram em espera longa: é preferível mostrar o erro ao usuário do que
bloquear a requisição. As evidências principais são os commits \commit{cb69179},
\commit{8eb4ae0}, \commit{6e6ab2d} e \commit{4196880}.

\subsection{Ingestão e indexação de documentos}

O baseline baixava PDFs de uma pasta do Drive e recriava o índice. O grupo ampliou
esse fluxo em cinco frentes.

\begin{enumerate}
  \item \textbf{Listagem recursiva e paginada.} O \arquivo{DriveService} percorre
        subpastas e continua a listagem além dos primeiros 100 itens, usando
        \texttt{nextPageToken}. Também valida se o ID informado aponta para uma
        pasta.
  \item \textbf{Novos formatos.} Além de PDFs, o sistema baixa planilhas
        \texttt{.xlsx} e exporta Google Docs como PDF. Arquivos \texttt{.xls} e
        Google Sheets nativos continuam fora do escopo.
  \item \textbf{Nomes seguros.} Os arquivos baixados recebem o sufixo
        \texttt{--<id>} para evitar colisões entre documentos com o mesmo título.
        Os IDs são removidos na apresentação das fontes ao usuário.
  \item \textbf{Indexação atômica.} O download ocorre em uma pasta temporária. Se
        a pasta estiver vazia ou um download falhar, a sincronização falha sem
        substituir o índice anterior.
  \item \textbf{Cache de embeddings.} O índice mantém um cache em disco dos
        vetores já calculados. Apenas os trechos novos são enviados ao servidor,
        e uma execução interrompida retoma de onde parou. O cache fica dentro da
        pasta do índice e é separado por modelo.
\end{enumerate}

Também foi criado o comando \texttt{python -m src.index\_documents}, que aceita um
caminho local ou a opção \texttt{--drive}. As evidências principais são
\commit{cb69179}, \commit{8eb4ae0}, \commit{952edd0} e \commit{6e6ab2d}.

\subsection{Fichas de alunos e consultas estruturadas}

Esta é a adição funcional mais extensa. O baseline não interpretava planilhas; o
grupo criou o módulo \arquivo{backend/src/services/spreadsheets.py} para ler
fichas de alunos e planilhas de acompanhamento.

O módulo reconhece diferentes tipos de aba:

\begin{itemize}
  \item \textbf{Perfil:} pares pergunta/resposta da ficha do aluno.
  \item \textbf{Alura:} cursos, carga horária, progresso, conclusão e avaliação.
  \item \textbf{Interações:} registros datados de mentores e coordenação.
  \item \textbf{Avaliações:} atividades, módulos, formações e comentários.
  \item \textbf{Scorecard:} semanas, horas estudadas, frequência, desempenho e
        legendas de nota.
\end{itemize}

Cada registro indexado recebe metadados como \texttt{kind}, \texttt{owner},
\texttt{sheet}, \texttt{cells}, \texttt{data}, \texttt{frequencia} e
\texttt{horas}. Isso permite que a busca diferencie um trecho genérico de uma
pergunta sobre a última interação de um aluno, por exemplo.

O fluxo de consulta foi ajustado para:

\begin{itemize}
  \item identificar o dono da ficha pelo campo de nome do aluno, com fallback para
        o nome do arquivo;
  \item filtrar linhas de planilhas coletivas antes de indexar;
  \item restringir a busca ao arquivo da pessoa citada;
  \item pedir desambiguação quando há mais de uma pessoa com o mesmo nome;
  \item não cair para a planilha inteira quando a identidade não é inequívoca;
  \item retornar citações com aba e intervalo de células.
\end{itemize}

As consultas suportadas incluem última interação, última nota do mentor, última
nota da coordenação, última semana com horas, resumo de cursos, cursos comentados,
semanas com falta e escolaridade. A ausência de um registro é tratada como
ausência de informação, não como resposta negativa inventada.

As evidências principais são \commit{e8b37b4}, \commit{2548cc4} e
\commit{fb61834}. O commit \commit{4376930}, ainda em branch separada, amplia as
consultas explícitas de cursos e faltas.

\subsection{Autenticação, autorização e segurança}

No baseline, o frontend obtinha um token Google, mas o backend não o validava. O
grupo passou a validar o ID token no backend com
\path{google.oauth2.id_token.verify_oauth2_token}, usando o Client ID
configurado e uma tolerância de 30 segundos para relógios dessincronizados.

Também foram adicionados:

\begin{itemize}
  \item \path{ALLOWED_EMAIL_DOMAIN}, com padrão
        \texttt{institutotadao-itt.org.br}; contas de outros domínios recebem 403;
  \item verificação de \path{email_verified};
  \item \path{AUTH_DEV_MODE} e o token \path{dev-token} apenas para
        desenvolvimento local;
  \item \path{ADMIN_EMAILS}, uma allowlist para o endpoint
        \path{/admin/sync-knowledge};
  \item mensagens de erro padronizadas em português para 401, 403, 429, 500 e 503;
  \item execução da sincronização em thread, para não bloquear a API;
  \item remoção do Firebase, que não era usado.
\end{itemize}

As evidências principais são \commit{39e20ce}, \commit{8e994e0},
\commit{b64ea6c}, \commit{574db76}, \commit{5b3f952} e \commit{6c7496e}.

\subsection{Chat, streaming, fontes e experiência de uso}

O endpoint \texttt{POST /chat/stream} transmite a resposta em NDJSON (Newline Delimited JSON). Cada linha é
um objeto JSON com um campo \texttt{type}:

\begin{lstlisting}[language=Python,caption={Formato simplificado dos eventos do endpoint de streaming.},label={lst:ndjson}]
{"type": "sources", "data": [...]}
{"type": "token", "data": "trecho da resposta"}
{"type": "done"}
\end{lstlisting}

As fontes são enviadas antes dos tokens, o que permite ao frontend preparar a
interface enquanto o texto chega. O botão ``Parar resposta'' usa
\texttt{AbortController}; respostas interrompidas ou com erro não são enviadas
como contexto nas perguntas seguintes. O histórico é mantido no
\texttt{sessionStorage}, separado por usuário, e limitado às últimas 12 mensagens.

A interface também ganhou o toggle ``Ver fontes'', que mostra arquivo, página, aba
e intervalo de células. O endpoint antigo \texttt{POST /chat/query} foi mantido
para compatibilidade e para o script de avaliação. As evidências principais são
\commit{303207e}, \commit{92e6ad2}, \commit{5b2c4fb}, \commit{4196880} e
\commit{c60f91c}.

\subsection{Docker e operação}

O grupo adicionou um \arquivo{compose.yaml} com dois serviços:

\begin{itemize}
  \item \textbf{backend:} imagem Python 3.12, healthcheck em \path{/health},
        volume nomeado para o índice FAISS e montagem somente leitura das
        credenciais;
  \item \textbf{frontend:} build multi-stage com Nginx, dependente do backend
        saudável.
\end{itemize}

As portas são publicadas apenas em \texttt{127.0.0.1}, coerente com o uso local e
com o modo de desenvolvimento. O guia \arquivo{RUN_CHATBOT.md} descreve a
instalação, a configuração do Drive, a indexação e a execução com Docker. Na raiz,
o comando \texttt{npm run dev} inicia backend e frontend com
\texttt{concurrently}. As evidências principais são \commit{952edd0},
\commit{18fde7e}, \commit{91b28a1}, \commit{c11f6d6} e \commit{0979185}.

\subsection{Avaliação de modelos}

O grupo criou o script \arquivo{backend/eval/run_evaluation.py}, que lê um CSV de
perguntas com \texttt{id}, \texttt{question}, \texttt{expected\_answer},
\texttt{expected\_source} e \texttt{has\_answer}, executa o fluxo do chatbot e
grava respostas, fontes retornadas, tempo e erro. O script permite retomar apenas
as perguntas com erro e sanitiza o nome do modelo no nome do arquivo CSV.

O commit \commit{4df4959}, da branch \path{vitor/avaliacao-modelos}, amplia essa
base com três scripts:

\begin{itemize}
  \item \arquivo{build_questions.py}: gera perguntas a partir de PDFs
        institucionais e de fichas, calculando respostas esperadas localmente;
  \item \arquivo{compare_models.py}: executa a avaliação em vários modelos e
        gera dois arquivos CSV, um de comparação e outro de resumo;
  \item \arquivo{prepare_review.py}: separa as perguntas de PDF para revisão
        manual e atribui notas automáticas às perguntas de planilha, sem expor
        dados de alunos.
\end{itemize}

As evidências principais são \commit{83b2c2e}, \commit{aed567e} e
\commit{4df4959}.

\subsection{Testes e documentação}

O baseline não tinha testes automatizados. O grupo criou quatro módulos, com 62
testes no total:

\begin{itemize}
  \item \arquivo{test_auth.py}: token ausente, inválido, domínio, modo dev,
        administradores, rate limit e erros;
  \item \arquivo{test_google_drive.py}: paginação, subpastas, exportação,
        nomes duplicados, planilhas, cache e falhas;
  \item \arquivo{test_openwebui.py}: histórico, triagem, RAG, fontes, retry 429,
        embeddings e erros;
  \item \arquivo{test_student_sheets.py}: identidade, ambiguidade, datas,
        células, faltas, escolaridade e compatibilidade.
\end{itemize}

A documentação foi consolidada no \arquivo{README.md}, no
\arquivo{RUN_CHATBOT.md}, no \arquivo{backend/README.md} e no
\arquivo{backend/DOCUMENTACAO_API.md}. As evidências são os commits de
documentação e os arquivos em \arquivo{backend/tests/}.

\subsection{Quadro-resumo das adições}

A Tabela~\ref{tab:antesdepois} resume a evolução entre o baseline e o estado atual.

\begin{table}[H]
\centering
\small
\begin{tabularx}{\textwidth}{@{}p{3.7cm} >{\raggedright\arraybackslash}X >{\raggedright\arraybackslash}X@{}}
\toprule
Área & Baseline (\commit{087a682}) & Estado atual (\commit{99da614}) \\
\midrule
Modelo de linguagem & Gemini & OpenWebUI/Unicamp com Gemma 3 12B; Gemini como fallback \\
Embeddings & Gemini & Nomic via OpenWebUI, com cache local \\
Documentos & PDF & PDF, Google Docs exportados, \texttt{.xlsx} e fichas de alunos \\
Autenticação & Apenas no frontend & Token validado no backend, domínio e allowlist de admin \\
Chat & Resposta JSON única & JSON e NDJSON com streaming de tokens \\
Fontes & Lista simples & Toggle com arquivo, página, aba e células \\
Operação & Docker apenas no backend & Docker Compose, Nginx, healthcheck e volume FAISS \\
Testes & Nenhum & 62 testes em quatro módulos \\
Avaliação & Nenhuma & Script de avaliação e comparação multi-modelo em branch \\
\bottomrule
\end{tabularx}
\caption{Quadro-resumo das adições em relação ao baseline.}
\label{tab:antesdepois}
\end{table}

% ======================================================================
\section{O que foi feito por cada participante}
% ======================================================================

\subsection{Critério de atribuição}

A descrição abaixo segue a autoria dos commits registrada no Git. Commits de merge
são contados separadamente, e branches não integradas são marcadas como tal. O
número de commits não deve ser lido como medida de esforço: um commit grande pode
concentrar uma funcionalidade inteira, enquanto vários commits pequenos podem
corrigir detalhes de uma mesma frente.

A Tabela~\ref{tab:contribuicoes} apresenta o consolidado.

\begin{table}[H]
\centering
\small
\begin{tabularx}{\textwidth}{@{}p{2.5cm} p{1.7cm} X p{2.4cm}@{}}
\toprule
Participante & Commits (all/main) & Frentes principais & Status \\
\midrule
Henrique & 28/28 & Autenticação, OpenWebUI, fontes, cache, UX,
avaliação inicial e documentação & \texttt{main} \\
Guilherme & 6/5 & OpenWebUI/Drive, fichas, escolaridade, merge do PR \#3 &
\texttt{main} + branch \\
Renilson & 7/7 & Docker/Compose, streaming, toggle de fontes & \texttt{main} \\
Vitor & 1/0 & Avaliação multi-modelo e revisão & branch não integrada \\
\bottomrule
\end{tabularx}
\caption{Consolidado de contribuições por participante.}
\label{tab:contribuicoes}
\end{table}

\subsection{Henrique Ribeiro Medeiros da Silva (243172)}

Henrique é o autor com maior número de commits no histórico: 28, todos alcançáveis
a partir da \texttt{main}. As frentes principais foram:

\begin{itemize}
  \item \textbf{Base de execução:} fixação do Python 3.12, inclusão do
        \texttt{slowapi}, \arquivo{.dockerignore}, regeneração de
        \arquivo{requirements.txt}, correção do \texttt{npm install} na raiz e
        remoção do Firebase sem uso.
  \item \textbf{Frontend:} leitura de \path{VITE_API_URL}, correção do
        timestamp das respostas do assistente, carregamento robusto do botão do
        Google, botão de parar resposta e mensagens de erro.
  \item \textbf{Autenticação e autorização:} validação do token Google no backend,
        restrição ao domínio do ITT, modo de desenvolvimento, allowlist de
        administradores e padronização de erros.
  \item \textbf{RAG e robustez:} fontes sob demanda, esclarecimento de nomes
        repetidos, cache de embeddings, busca restrita à planilha da pessoa,
        retry para 429 na indexação e processamento de perguntas em thread.
  \item \textbf{Avaliação e documentação:} script inicial de avaliação,
        correção do nome do CSV, README unificado, \arquivo{RUN_CHATBOT.md} e
        documentação da API.
\end{itemize}

Commits representativos: \commit{5b3f952}, \commit{39e20ce}, \commit{8eb4ae0},
\commit{b64ea6c}, \commit{8e994e0}, \commit{6e6ab2d}, \commit{303207e},
\commit{92e6ad2}, \commit{134fe20} e \commit{83b2c2e}. Henrique também assinou os
merges dos PRs \#1 e \#2.

\subsection{Guilherme de Godoi Monteiro (237031)}

Guilherme tem seis commits autorais: cinco no \texttt{main} e um na branch
\texttt{fix/\allowbreak consultas-\allowbreak alura-\allowbreak faltas}. As frentes principais foram:

\begin{itemize}
  \item \textbf{Integração inicial com a Unicamp:} o commit \commit{cb69179}
        introduz o provedor OpenWebUI, reestrutura o acesso ao Google Drive, cria
        o CLI de indexação, adiciona os primeiros testes de Drive/OpenWebUI e
        atualiza a documentação de execução. É um commit grande, com 23 arquivos
        e cerca de 1.287 inserções.
  \item \textbf{Fichas de alunos:} os commits \commit{e8b37b4}, \commit{2548cc4}
        e \commit{fb61834} implementam consultas por registro e intervalo de
        células, recuperação da escolaridade, isolamento entre fichas e
        compatibilidade com índices anteriores.
  \item \textbf{Integração:} \commit{99da614} é o merge do PR \#3, que traz a
        leitura de fichas para a \texttt{main}.
  \item \textbf{Branch adicional:} \commit{4376930} responde a consultas
        explícitas de cursos e faltas com dados da ficha, mas ainda não foi
        integrado.
\end{itemize}

\subsection{Renilson Conceição de Luna Júnior (247358)}

Renilson tem sete commits, todos no \texttt{main}. As frentes principais foram:

\begin{itemize}
  \item \textbf{Teste de modelo:} o documento \arquivo{teste/teste-modelo.md}
        registra o planejamento inicial de perguntas e um script de avaliação
        com Gemini.
  \item \textbf{Docker e operação:} o commit \commit{952edd0} adiciona o
        \arquivo{compose.yaml}, o Dockerfile do frontend, a configuração do Nginx,
        o endpoint \path{/health} e instruções no \arquivo{RUN_CHATBOT.md}.
  \item \textbf{Streaming:} os commits \commit{5b2c4fb} e \commit{4196880}
        implementam o streaming de texto no backend e no frontend, incluindo o
        toggle de fontes e o consumo de eventos SSE (Server-Sent Events) do OpenWebUI.
  \item \textbf{Integração:} \commit{b64a76e}, \commit{496789d} e
        \commit{c60f91c} fazem merge e ajustes de compatibilidade entre a branch
        de streaming e a \texttt{main}.
\end{itemize}

\subsection{Vitor Ferreira dos Santos de Lacerda (231524)}

Vitor tem um commit, \commit{4df4959}, na branch
\path{vitor/avaliacao-modelos}, ainda não integrada à \texttt{main}. O commit
adiciona 401 linhas em três scripts:

\begin{itemize}
  \item \arquivo{build_questions.py}: gera o CSV de perguntas, combinando
        perguntas fixas sobre PDFs institucionais com perguntas derivadas das
        fichas; as respostas esperadas são calculadas localmente.
  \item \arquivo{compare_models.py}: executa a avaliação em vários modelos,
        compara respostas, fontes e tempos e gera arquivos de resumo.
  \item \arquivo{prepare_review.py}: prepara a revisão manual das perguntas de
        PDF e atribui notas automáticas às perguntas de planilha, preservando a
        privacidade dos dados de alunos.
\end{itemize}

No estado atual, a branch está 1 commit à frente e 12 atrás da \texttt{main}. Para
ser apresentada como entrega consolidada, precisa de rebase ou merge e de uma
revisão dos conflitos com os scripts de avaliação já existentes.

% ======================================================================
\section{Considerações finais}
% ======================================================================

O projeto evoluiu de um chatbot RAG experimental, limitado a PDFs e ao provedor
Gemini, para uma aplicação interna com autenticação, ingestão de documentos do
Drive, consultas estruturadas a fichas de alunos, streaming, fontes detalhadas,
testes automatizados e execução reprodutível com Docker Compose.

As adições mais relevantes são a validação de identidade no backend, a leitura
estruturada de planilhas e fichas, o provedor OpenWebUI da Unicamp e o streaming
com possibilidade de interrupção. Essas mudanças aproximam o sistema de um uso real
pela equipe do ITT, mas não eliminam todas as limitações.

Os principais trabalhos futuros são:

\begin{itemize}
  \item implementar permissões por perfil, para que cada usuário veja apenas os
        documentos e alunos aos quais tem direito;
  \item adicionar OCR para PDFs digitalizados;
  \item suportar Google Sheets nativos e arquivos \texttt{.xls};
  \item criar uma política de expiração e limpeza do cache de embeddings;
  \item integrar a avaliação multi-modelo à \texttt{main} e automatizar a revisão;
  \item adicionar integração contínua para executar os testes a cada alteração;
  \item monitorar custo, latência e taxa de erro do provedor de LLM em produção.
\end{itemize}

% ======================================================================
% Referências
% ======================================================================
\phantomsection
\addcontentsline{toc}{section}{Referências}
\begin{thebibliography}{9}

\bibitem{repo}
HENRIQUE SILVA et al.
\textbf{chatbot-itt-2026-2}.
GitHub, 2026.
Disponível em: \url{https://github.com/henriquermdsilva/chatbot-itt-2026-2}.
Acesso em: 27 nov. 2026.

\bibitem{antecessor}
CHRIS ARAÚJO et al.
\textbf{chatbot-itt-2026-1}.
GitHub, 2026.
Disponível em: \url{https://github.com/chraujo/chatbot-itt-2026-1}.
Acesso em: 27 nov. 2026.

\bibitem{programacao}
PROGRAMAÇÃO DEV.
\textbf{chatbot-itt-}.
GitHub, 2025.
Disponível em: \url{https://github.com/programacao-dev/chatbot-itt-}.
Acesso em: 27 nov. 2026.

\bibitem{fastapi}
TIANGOLO, Sebastián Ramírez.
\textbf{FastAPI documentation}.
Disponível em: \url{https://fastapi.tiangolo.com/}.
Acesso em: 27 nov. 2026.

\bibitem{langchain}
LANGCHAIN.
\textbf{LangChain documentation}.
Disponível em: \url{https://python.langchain.com/}.
Acesso em: 27 nov. 2026.

\bibitem{langgraph}
LANGCHAIN.
\textbf{LangGraph documentation}.
Disponível em: \url{https://langchain-ai.github.io/langgraph/}.
Acesso em: 27 nov. 2026.

\bibitem{faiss}
META AI.
\textbf{FAISS: A library for efficient similarity search}.
Disponível em: \url{https://github.com/facebookresearch/faiss}.
Acesso em: 27 nov. 2026.

\bibitem{openwebui}
OPEN WEBUI.
\textbf{Open WebUI documentation}.
Disponível em: \url{https://docs.openwebui.com/}.
Acesso em: 27 nov. 2026.

\bibitem{drive}
GOOGLE.
\textbf{Google Drive API documentation}.
Disponível em: \url{https://developers.google.com/workspace/drive/api}.
Acesso em: 27 nov. 2026.

\end{thebibliography}

% ======================================================================
% Apêndices
% ======================================================================
\end{document}
